\section*{الفصل \ref{ch:g11:hormones-and-reproduction} --- الضبط الهرموني للتكاثر}

\begin{solution}{exo:g11:hormones-and-reproduction:1}
\omterm{def:g11:hormones-and-reproduction:hormone}{الهرمون} جزيء تفرزه غدة في الدم ويؤثر في \omterm{def:g10:cells-common-unit:cell}{خلايا} بعيدة تحمل مستقبله.
\omterm{def:g11:hormones-and-reproduction:hormone}{والتغذية الراجعة} هي أثر الناتج النهائي في الغدد التي في أول السلسلة:
سالبة حين يخفض الأمر (فيبقي مستوى ثابتًا)، وموجبة حين يزيده (فينتج
اندفاعة).
\end{solution}

\begin{solution}{exo:g11:hormones-and-reproduction:2}
الوطاء (GnRH)، والنخامية (LH وFSH)، والقند (التستوستيرون في الخصية؛
والإستروجينات والبروجسترون في المبيض).
\end{solution}

\begin{solution}{exo:g11:hormones-and-reproduction:3}
الأنيبيبات المنوية تنتج النطاف؛ \omterm{def:g10:cells-common-unit:cell}{والخلايا} الواقعة بينها تفرز
التستوستيرون.
\end{solution}

\begin{solution}{exo:g11:hormones-and-reproduction:4}
الطور الجريبي (الأيام 1--13): جريب نامٍ، وإستروجينات. والإباضة (اليوم
14): اندفاعة LH. والطور الأصفري (الأيام 15--28): جسم أصفر،
وبروجسترون وإستروجينات.
\end{solution}

\begin{solution}{exo:g11:hormones-and-reproduction:5}
تنكس الجسم الأصفر في نهاية الدورة: فهبوط البروجسترون والإستروجينات
يجعل بطانة الرحم الغليظة تتفكك وتنزف.
\end{solution}

\begin{solution}{exo:g11:hormones-and-reproduction:6}
ذروة الإستروجين في اليوم 12، وذروة LH في اليوم 13، والإباضة في اليوم
14، وذروة البروجسترون في اليوم 21.
\end{solution}

\begin{solution}{exo:g11:hormones-and-reproduction:7}
الخصاء يزيل التستوستيرون، ومنه تثبيطه للوطاء والنخامية: فيرتفع LH.
والتستوستيرون المحقون يعيد التثبيط فيهبط LH --- أي \omterm{def:g11:hormones-and-reproduction:hormone}{تغذية راجعة} سالبة
في الاتجاهين.
\end{solution}

\begin{solution}{exo:g11:hormones-and-reproduction:8}
في اليوم 12 يكون مستوى الإستروجين مرتفعًا ومستمرًا، وهذا يقلب \omterm{def:g11:hormones-and-reproduction:hormone}{التغذية
الراجعة} إلى موجبة ويطلق الاندفاعة. أما في اليوم 20 فالإستروجينات
معتدلة والبروجسترون مرتفع، \omterm{def:g11:hormones-and-reproduction:hormone}{والتغذية الراجعة} سالبة: فيبقى LH منخفضًا.
\end{solution}

\begin{solution}{exo:g11:hormones-and-reproduction:9}
اليوم 21: فالطور الأصفري يدوم نحو 14 يومًا أيًّا كان طول الدورة،
فالإباضة قبل الحيض التالي بأربعة عشر يومًا؛ والأسبوع الزائد في الطور
الجريبي.
\end{solution}

\begin{solution}{exo:g11:hormones-and-reproduction:10}
إستروجين الحبة الثابت المعتدل وهرمونها الشبيه بالبروجسترون يبقيان
\omterm{def:g11:hormones-and-reproduction:hormone}{التغذية الراجعة} سالبة طوال الشهر: فيبقى FSH أدنى من أن ينضج جريب، فلا
تقع ذروة إستروجين، ولا \omterm{def:g11:hormones-and-reproduction:hormone}{تغذية راجعة} موجبة، ولا اندفاعة LH، ولا إباضة.
\end{solution}

\begin{solution}{exo:g11:hormones-and-reproduction:11}
التستوستيرون الخارجي يثبّط وطاءه ونخاميته؛ فيهبط LH وFSH؛ وتكف خصيتاه،
وقد كفّ التنبيه عنهما، عن الإفراز وعن صنع النطاف، فتضمران.
\end{solution}

\begin{solution}{exo:g11:hormones-and-reproduction:12}
من دون تثبيط الإستروجين، يرتفع GnRH وFSH؛ فتنضج الجريبات؛ وذروة
إستروجينها (وهي تؤثر في النخامية، وهي لا تزال قادرة على الاستجابة)
تطلق اندفاعة LH والإباضة. ويستعمل حين تخفق الإباضة لأن الجريبات تتلقى
FSH قليلًا جدًا.
\end{solution}

\begin{solution}{exo:g11:hormones-and-reproduction:13}
الجريب ينمو ويفرز (فالمبيض وFSH يعملان)، لكن ذروة الإستروجين لا تنتج
اندفاعة: \omterm{def:g11:hormones-and-reproduction:hormone}{فالتغذية الراجعة} الموجبة تخفق في الوطاء أو في النخامية. وحقنة
من \omterm{def:g11:hormones-and-reproduction:hormone}{هرمون} شبيه \omterm{def:g11:hormones-and-reproduction:hormone}{بهرمون} LH في وقت الذروة تستطيع إطلاق الإباضة.
\end{solution}

\begin{solution}{exo:g11:hormones-and-reproduction:14}
إذا أبقي الجسم الأصفر حيًا واصل إفراز البروجسترون: فتحفظ البطانة (فلا
طمث) وتبقي \omterm{def:g11:hormones-and-reproduction:hormone}{التغذية الراجعة} السالبة هرموني LH وFSH
منخفضين (فلا جريب جديد ولا إباضة). \omterm{def:g11:hormones-and-reproduction:hormone}{والهرمون} يظهر في الدم والبول في غضون أيام من الانغراس، وهو ما
تكشفه الاختبارات.
\end{solution}

\begin{solution}{exo:g11:hormones-and-reproduction:15}
$12 \times 0.6 \times 0.4 \approx 2.9$: أي نحو 3 أجنة. واحتمال ألا
تحدث ولادة في ثلاث نقلات: $0.65^3 \approx 0.27$؛ واحتمال ولادة واحدة
على الأقل: نحو 73\%.
\end{solution}

\begin{solution}{pb:g11:hormones-and-reproduction:1}
\textbf{1.} اليوم 14: فاندفاعة LH في اليوم 13 تسبق الإباضة بيوم،
والبروجسترون، ولا يظهر إلا بعد الإباضة، يرتفع من اليوم 14.

\textbf{2.} الطور الجريبي الأيام 1--13 (13 يومًا)؛ والطور الأصفري
الأيام 14--27 أو 28 (نحو 14 يومًا).

\textbf{3.} الجريب النامي في اليوم 10؛ والجسم الأصفر في اليوم 22،
للإستروجينات وللبروجسترون معًا.

\textbf{4.} اليوم 6: رقيقة، بدأت تغلظ تحت الإستروجينات. واليوم 22:
غليظة مفرزة تحت البروجسترون، مهيأة لجنين. واليوم 28: تتفكك وتنزف، وقد
هبط الهرمونان.

\textbf{5.} في البداية تكون \omterm{def:g11:hormones-and-reproduction:hormone}{هرمونات} الدورة السابقة قد هبطت وهبط
تثبيطها معها، فيرتفع FSH؛ وفي الطور الأصفري يثبّطه البروجسترون
والإستروجينات تثبيطًا قويًا.

\textbf{6.} \omterm{def:g11:hormones-and-reproduction:hormone}{تغذية راجعة} سالبة: فإستروجينات الجريب المرتفعة تثبّط
النخامية، فيهبط FSH.

\textbf{7.} \omterm{def:g11:hormones-and-reproduction:hormone}{تغذية راجعة} موجبة: فالإستروجينات عند أعظميها نحو يومين
تنبّه الوطاء والنخامية بدل أن تثبّطهما، فيندفع LH.

\textbf{8.} سالبة: فالبروجسترون مع إستروجينات معتدلة يثبّط النخامية؛
ومن دون ذروة إستروجين مستمرة، لا \omterm{def:g11:hormones-and-reproduction:hormone}{تغذية راجعة} موجبة ولا اندفاعة ثانية.

\textbf{9.} ينهار البروجسترون والإستروجينات؛ فتتفكك البطانة وتنزف في
اليوم الأول من الدورة التالية؛ ويرتفع FSH، وقد تحرر من التثبيط.

\textbf{10.} هبوط هرموني الطور الأصفري يزيل \omterm{def:g11:hormones-and-reproduction:hormone}{التغذية الراجعة} السالبة،
فيرتفع FSH، ويبدأ جريب جديد بالنمو: فنهاية دورة إشارة للتي تليها.

\textbf{11.} FSH وLH منخفضان طوال الشهر (\omterm{def:g11:hormones-and-reproduction:hormone}{فالتغذية الراجعة} السالبة
تحفظها \omterm{def:g11:hormones-and-reproduction:hormone}{هرمونات} الحبة)؛ ولا ينضج جريب.

\textbf{12.} لا ذروة (فلا جريب ناضج)، ولا اندفاعة، ولا إباضة.

\textbf{13.} تهبط \omterm{def:g11:hormones-and-reproduction:hormone}{هرمونات} الحبة وتتفكك البطانة، وهي رقيقة، وتنزف؛ وهو
نزف انسحاب، لا نهاية طور أصفري --- إذ لم يكن ثمة جسم أصفر.

\textbf{14.} من دون \omterm{def:g11:hormones-and-reproduction:hormone}{هرمونات} الحبة ترتفع \omterm{def:g11:hormones-and-reproduction:hormone}{التغذية الراجعة}، فيرتفع FSH،
وقد يبدأ جريب بالنمو والإفراز؛ فإذا بلغ ذروة إستروجينه، أمكن أن تتلوه
اندفاعة LH وإباضة، ولم تعد البطانة الرقيقة والمخاط الغليظ محفوظين.

\textbf{15.} البروجسترون وحده يثبّط الوطاء والنخامية بما يكفي لمنع
اندفاعة LH، ويغلّظ مخاط عنق الرحم.

\textbf{16.} الوطاء أو النخامية: فالمبيض غير منبَّه. ونبضات من GnRH
(إن كانت النخامية تعمل) أو حقن FSH وLH تعيد نمو الجريبات والإباضة،
متجاوزةً الطبقة المفقودة.

\textbf{17.} الإخصاب في المختبر: بييضات تجمع من المبيض، وتخصَّب في
طبق، ويوضع الجنين في الرحم --- أي لقاء الأمشاج والنقل الذي كانت قناة
البيض لتقوم به.

\textbf{18.} اندفاعة LH الطبيعية؛ فالإباضة تتلو الاندفاعة بنحو 36
ساعة، فتجمع البييضات قبيل أن تطلق وتضيع.

\textbf{19.} طبيعيًا، تخفض إستروجينات الجريب الرائد المرتفعة FSH،
فتنكس الجريبات الأصغر، وقد حرمت FSH؛ ولا ينجو إلا الأكبر. أما FSH
المحقون فيتجاوز هذا الانتقاء ويبقيها كلها نامية.

\textbf{20.} الإباضة في اليوم 14؛ \omterm{def:g11:hormones-and-reproduction:hormone}{والتغذية الراجعة} السالبة التي خفضت
FSH مع نمو الجريب، ثم \omterm{def:g11:hormones-and-reproduction:hormone}{التغذية الراجعة} الموجبة التي حوّلت ذروة
الإستروجين إلى اندفاعة LH؛ والحبة أبقت \omterm{def:g11:hormones-and-reproduction:hormone}{التغذية الراجعة} سالبة طوال
الشهر، فلم ينضج جريب ولم تقع اندفاعة.
\end{solution}
